\documentclass[11pt,a4paper]{article}

\usepackage[margin=1in]{geometry}
\usepackage{amsmath,amssymb}
\usepackage{booktabs}
\usepackage{graphicx}
\usepackage{hyperref}
\usepackage{caption}
\usepackage{subcaption}
\usepackage{array}
\usepackage{xcolor}

\hypersetup{colorlinks=true, linkcolor=blue, citecolor=blue, urlcolor=blue}

\newcommand{\ncsim}{\texttt{ncsim}}
\newcommand{\MB}{\,\text{MB}}
\newcommand{\MBs}{\,\text{MB/s}}
\newcommand{\Mbps}{\,\text{Mbps}}
\newcommand{\dBm}{\,\text{dBm}}
\newcommand{\dB}{\,\text{dB}}
\newcommand{\GHz}{\,\text{GHz}}
\newcommand{\MHz}{\,\text{MHz}}

\title{Verification of the CSMA/CA Bianchi Interference Model\\in the \ncsim{} Networked Compute Simulator}
\author{Autonomous Networks Research Group (ANRG)\\University of Southern California}
\date{March 2026}

\begin{document}
\maketitle

%% ────────────────────────────────────────────────────────────────
\begin{abstract}
We present a systematic verification of the \texttt{csma\_bianchi} wireless interference model implemented in \ncsim{}, a headless discrete-event simulator for networked computing research. The model separates two distinct 802.11 interference mechanisms: (1)~CSMA/CA contention-domain time-sharing modeled by Bianchi's saturation throughput analysis, and (2)~hidden-terminal SINR degradation from links outside the carrier-sensing range. Nine experiments compare simulator output against closed-form analytical predictions across a range of topologies, from single-link baselines to five-link cascading hidden-terminal scenarios and combined bandwidth-sharing--plus--interference configurations. All~108 individual test points match predictions within a 1\% tolerance, validating the correctness of the interference model, the execution engine's dynamic bandwidth recalculation, and the per-link fair-sharing mechanism.
\end{abstract}

%% ────────────────────────────────────────────────────────────────
\section{Introduction}
\label{sec:intro}

Accurate modeling of wireless link interference is essential for any networked-computing simulator that aims to produce realistic transfer-time estimates. In multi-hop or multi-link wireless networks, concurrent transmissions interact through two fundamentally different mechanisms:

\begin{enumerate}    \item \textbf{Contention-domain time-sharing.} Links whose transmitters can sense each other (i.e., within carrier-sensing range) contend for channel access via the 802.11 CSMA/CA protocol. These links share airtime but do not transmit simultaneously, so they do not cause signal-level interference at each other's receivers.
    \item \textbf{Hidden-terminal SINR degradation.} Links whose transmitters \emph{cannot} sense each other may transmit simultaneously, causing co-channel interference at the receiver. The resulting signal-to-interference-plus-noise ratio (SINR) determines the achievable MCS and data rate.
\end{enumerate}

The \ncsim{} simulator implements these two mechanisms in its \texttt{csma\_bianchi} interference model. The combined effective data rate for a link is:
\begin{equation}
    R_{\text{eff}} = \frac{R_{\text{base}} \cdot f_{\text{SINR}} \cdot \eta(n)}{n \cdot N_{\text{flows}}}
    \label{eq:combined}
\end{equation}
where $R_{\text{base}}$ is the SNR-only PHY rate, $f_{\text{SINR}}$ captures hidden-terminal degradation, $\eta(n)/n$ is the Bianchi per-station throughput share among $n$ contending links, and $N_{\text{flows}}$ is the number of concurrent data transfers sharing the same physical link.

This report describes nine experiments that verify Equation~\eqref{eq:combined} and the simulator's dynamic recalculation behavior against closed-form predictions. Section~\ref{sec:model} describes the interference model in detail. Sections~\ref{sec:exp1}--\ref{sec:exp9} present each experiment. Section~\ref{sec:conclusions} summarizes the results.

%% ────────────────────────────────────────────────────────────────
\section{System Model}
\label{sec:model}

\subsection{RF Parameters}

All experiments use the default \texttt{RFConfig} parameters listed in Table~\ref{tab:rf}.

\begin{table}[h]
\centering
\caption{Default RF configuration.}
\label{tab:rf}
\begin{tabular}{lrl}
\toprule
Parameter & Value & Unit \\
\midrule
Transmit power ($P_{\text{tx}}$) & 20.0 & dBm \\
Carrier frequency ($f$) & 5.0 & GHz \\
Path-loss exponent ($n_{\text{PL}}$) & 3.0 & -- \\
Noise floor ($N_0$) & $-95.0$ & dBm \\
CCA threshold & $-82.0$ & dBm \\
Channel width & 20 & MHz \\
WiFi standard & 802.11ax & -- \\
\bottomrule
\end{tabular}
\end{table}

\subsection{Path Loss and Rate Selection}

The log-distance path-loss model gives the received power at distance~$d$:
\begin{equation}
    P_{\text{rx}}(d) = P_{\text{tx}} - \text{PL}(d_0) - 10\, n_{\text{PL}} \log_{10}(d/d_0)
\end{equation}
where $\text{PL}(d_0) = 20\log_{10}(4\pi d_0 f/c)$ is the Friis free-space loss at reference distance $d_0 = 1\,$m. The signal-to-noise ratio is $\text{SNR} = P_{\text{rx}} - N_0$ (in dB). The highest 802.11ax MCS whose minimum SNR requirement is satisfied determines the PHY data rate.

\subsection{Carrier Sensing and Conflict Graph}

The carrier-sensing range is the maximum distance at which a transmission triggers the CCA:
\begin{equation}
    d_{\text{CS}} = d_0 \cdot 10^{(P_{\text{tx}} - \text{CCA} - \text{PL}(d_0)) / (10 \, n_{\text{PL}})}
\end{equation}
With the default parameters, $d_{\text{CS}} = 71.2\,$m. Two links \emph{conflict} (without RTS/CTS) if the transmitter of either link can sense any node of the other link. The set of all pairwise conflicts forms the \emph{conflict graph}.

\subsection{Bianchi MAC Efficiency}

For $n$ stations contending under 802.11 DCF, Bianchi's saturation-throughput analysis~\cite{bianchi2000} yields a MAC efficiency $\eta(n) \in (0,1]$ representing the fraction of channel time occupied by successful payload delivery. Each of the $n$ contending links obtains a throughput share of $\eta(n)/n$. For $n=1$, $\eta(1)=1.0$ (no contention overhead).

\subsection{Hidden-Terminal SINR}

Active links \emph{not} in the conflict graph of a given link are its hidden terminals. They may transmit simultaneously, so their aggregate interference is included in the SINR calculation:
\begin{equation}
    \text{SINR} = \frac{P_{\text{signal}}}{N_0 + \sum_{j \in \mathcal{H}} P_j}
\end{equation}
where $\mathcal{H}$ is the set of active hidden-terminal links and $P_j$ is the received interference power from hidden terminal~$j$'s transmitter at the victim receiver. A new MCS is selected based on the resulting SINR, yielding a rate $R_{\text{SINR}}$. The SINR factor is $f_{\text{SINR}} = R_{\text{SINR}} / R_{\text{base}}$.

\subsection{Combined Interference Factor}

The \texttt{csma\_bianchi} model composes the two mechanisms multiplicatively:
\begin{equation}
    f = f_{\text{SINR}} \cdot \frac{\eta(n)}{n}, \qquad f \in [0.01, \; 1.0]
\end{equation}
The effective per-flow bandwidth on a link carrying $N_{\text{flows}}$ concurrent transfers is then $R_{\text{base}} \cdot f / N_{\text{flows}}$.

\subsection{Dynamic Recalculation}

When any transfer starts or completes, \emph{all} other active links are recalculated. Partial progress is tracked via a \texttt{data\_remaining} field:
\begin{equation}
    D_{\text{rem}}' = D_{\text{rem}} - R_{\text{eff}} \cdot \Delta t
\end{equation}
where $\Delta t$ is the elapsed time since the last recalculation. A new completion time is then scheduled based on $D_{\text{rem}}'$ and the updated effective rate.

\subsection{Verification Methodology}

Each experiment generates a scenario YAML, runs it through \ncsim{}, and compares simulated transfer durations (or effective rates) against closed-form analytical predictions computed from the same RF functions. A match is declared when the relative error is within 1\%.

%% ────────────────────────────────────────────────────────────────
\section{Experiment~1: Link Length vs.\ Data Rate}
\label{sec:exp1}

\subsection{Setup}
A single 802.11ax link connects two nodes at varying distances $d \in \{1, 3, 5, 8, \ldots, 140\}\,$m. A single 10\,MB transfer is sent across the link with no other active links (no interference). The predicted rate is simply the SNR-only PHY rate $R_{\text{base}}(d)$.

\subsection{Purpose}
Validates the baseline PHY-layer pipeline: distance $\to$ path loss $\to$ SNR $\to$ MCS selection $\to$ data rate. This is a prerequisite for all subsequent experiments.

\subsection{Results}

\begin{table}[h]
\centering
\caption{Experiment~1: Single-link PHY rate vs.\ distance.}
\label{tab:exp1}
\begin{tabular}{rrrrrrr}
\toprule
$d$ (m) & PL (dB) & SNR (dB) & MCS & Predicted (MB/s) & Simulated (MB/s) & Match \\
\midrule
  1 & 46.42 & 68.58 & 11 & 17.9250 & 17.9250 & \checkmark \\
  5 & 67.39 & 47.61 & 11 & 17.9250 & 17.9250 & \checkmark \\
 12 & 78.80 & 36.20 &  9 & 14.3375 & 14.3375 & \checkmark \\
 20 & 85.45 & 29.55 &  7 & 10.7500 & 10.7500 & \checkmark \\
 30 & 90.73 & 24.27 &  5 &  8.6000 &  8.6000 & \checkmark \\
 50 & 97.39 & 17.61 &  3 &  4.3000 &  4.3000 & \checkmark \\
 75 &102.67 & 12.33 &  2 &  3.2250 &  3.2250 & \checkmark \\
 95 &105.75 &  9.25 &  1 &  2.1500 &  2.1500 & \checkmark \\
120 &108.80 &  6.20 &  0 &  1.0750 &  1.0750 & \checkmark \\
140 &110.81 &  4.19 & n/a &  0.0000 &  0.0000 & \checkmark \\
\bottomrule
\end{tabular}
\end{table}

All 20 test points match exactly. The rate decreases in discrete steps corresponding to MCS thresholds, and drops to zero beyond $d=130\,$m where the SNR falls below the minimum MCS~0 requirement of 5\,dB.

%% ────────────────────────────────────────────────────────────────
\section{Experiment~2: Parallel Link Separation}
\label{sec:exp2}

\subsection{Setup}
Two parallel horizontal links of length 30\,m are placed at vertical separation $s \in \{5, 10, \ldots, 200\}\,$m. Link~A runs from $n_0(0,0)$ to $n_1(30,0)$; Link~B from $n_2(0,s)$ to $n_3(30,s)$. Each carries a 10\,MB transfer. When $s \leq d_{\text{CS}} = 71.2\,$m, both links are in the conflict graph (Bianchi contention). When $s > d_{\text{CS}}$, they are hidden terminals (SINR degradation).

\subsection{Predictions}
\begin{itemize}    \item \textbf{Conflict regime} ($s \leq 71.2\,$m): $R = R_{\text{base}} \cdot \eta(2)/2 = 8.6 \times 0.4404 = 3.7871\MBs$.
    \item \textbf{Hidden-terminal regime} ($s > 71.2\,$m): The interferer distance is $d_I = \sqrt{30^2 + s^2}$. Both links see the same interferer geometry, so rates are symmetric: $R = R_{\text{SINR}}(30, d_I)$.
\end{itemize}

\subsection{Results}

\begin{table}[h]
\centering
\caption{Experiment~2: Two parallel links at varying separation.}
\label{tab:exp2}
\begin{tabular}{rcrrrr}
\toprule
$s$ (m) & In Conflict & Pred.\ $R_A$ & Sim.\ $R_A$ & Pred.\ $R_B$ & Sim.\ $R_B$ \\
\midrule
  5 & Yes & 3.7871 & 3.7871 & 3.7871 & 3.7871 \\
 20 & Yes & 3.7871 & 3.7871 & 3.7871 & 3.7871 \\
 50 & Yes & 3.7871 & 3.7871 & 3.7871 & 3.7871 \\
 70 & Yes & 3.7871 & 3.7871 & 3.7871 & 3.7871 \\
 75 & No  & 3.2250 & 3.2250 & 3.2250 & 3.2250 \\
 80 & No  & 3.2250 & 3.2250 & 3.2250 & 3.2250 \\
100 & No  & 4.3000 & 4.3000 & 4.3000 & 4.3000 \\
140 & No  & 6.4500 & 6.4500 & 6.4500 & 6.4500 \\
200 & No  & 6.4500 & 6.4500 & 6.4500 & 6.4500 \\
\bottomrule
\end{tabular}
\end{table}

All 15 test points match. In the conflict regime the rate is constant at $\eta(2)/2 \approx 0.44$ of the base rate regardless of separation. In the hidden-terminal regime, the rate initially drops (at $s=75\,$m the SINR degrades the MCS) then recovers as the interferer becomes more distant. Both links always receive the same rate, confirming symmetric behavior.

%% ────────────────────────────────────────────────────────────────
\section{Experiment~3: Two Transmitters to Same Receiver}
\label{sec:exp3}

\subsection{Setup}
Two transmitter nodes $n_1(d,0)$ and $n_2(0,d)$ both send 10\,MB to a shared receiver $n_0(0,0)$, positioned at 90$^\circ$ apart at equal distance~$d$. This tests the model when two links share a receiver, for $d \in \{5, 10, \ldots, 130\}\,$m.

\subsection{Predictions}
In the conflict regime ($d \leq 71.2\,$m, since each transmitter can sense $n_0$ at distance~$d$), the rate is $R_{\text{base}}(d) \cdot \eta(2)/2$. In the hidden-terminal regime, both interferers are at distance~$d$ from the shared receiver, causing equal-power interference. At $d=75\,$m this yields a SINR so low that the model clamps to the minimum factor of 0.01.

\subsection{Results}

\begin{table}[h]
\centering
\caption{Experiment~3: Two transmitters to shared receiver.}
\label{tab:exp3}
\begin{tabular}{rcrrrr}
\toprule
$d$ (m) & In Conflict & Pred.\ $R_{l10}$ & Sim.\ $R_{l10}$ & Pred.\ $R_{l20}$ & Sim.\ $R_{l20}$ \\
\midrule
  5 & Yes & 7.8934 & 7.8934 & 7.8934 & 7.8934 \\
 20 & Yes & 4.7338 & 4.7338 & 4.7338 & 4.7338 \\
 30 & Yes & 3.7871 & 3.7871 & 3.7871 & 3.7871 \\
 50 & Yes & 1.8935 & 1.8935 & 1.8935 & 1.8935 \\
 70 & Yes & 1.4201 & 1.4201 & 1.4201 & 1.4201 \\
 75 &  No & 0.0323 & 0.0323 & 0.0323 & 0.0323 \\
100 &  No & 0.0215 & 0.0215 & 0.0215 & 0.0215 \\
130 &  No & 0.0108 & 0.0108 & 0.0108 & 0.0108 \\
\bottomrule
\end{tabular}
\end{table}

All 18 test points match. The dramatic rate collapse at $d=75\,$m occurs because equal-power hidden-terminal interference at the shared receiver yields an extremely low SINR, demonstrating the model's minimum-factor clamping behavior.

%% ────────────────────────────────────────────────────────────────
\section{Experiment~4: Three Parallel Links}
\label{sec:exp4}

\subsection{Setup}
Three parallel 30\,m links are placed at $y = 0$, $y = s$, and $y = 2s$, labeled A (outer), B (middle), and C (outer). Each carries a 10\,MB transfer. As separation~$s$ increases, three interference regimes emerge:

\begin{itemize}    \item \textbf{All-conflict} ($s \leq 35.6\,$m): All three pairs are within $d_{\text{CS}}$. Each link gets $R_{\text{base}} \cdot \eta(3)/3$.
    \item \textbf{Mixed} ($35.6 < s \leq 71.2\,$m): Adjacent pairs (A--B, B--C) are in conflict, but the outer pair (A--C, at distance $2s$) are hidden terminals. A and C each contend with B ($\eta(2)/2$) and suffer SINR from the far hidden terminal. B contends with both ($\eta(3)/3$) but has no hidden terminals.
    \item \textbf{All-hidden} ($s > 71.2\,$m): No conflict-graph edges. Each link sees the others purely as hidden terminals with multi-interferer SINR.
\end{itemize}

\subsection{Multi-Phase Prediction}
Because A/C and B generally have different rates, one group finishes before the other. When links complete, the remaining links' interference conditions change, creating a second phase with updated rates. The analytical prediction accounts for this by computing:
\begin{enumerate}
    \item Phase~1 rates from the initial interference conditions.
    \item Which group finishes first (smallest $D/R$).
    \item Phase~2 rates for remaining links with reduced interference.
    \item Effective rate = $D / (\text{Phase~1 duration} + \text{Phase~2 duration})$.
\end{enumerate}

\subsection{Results}

\begin{table}[h]
\centering
\caption{Experiment~4: Three parallel links at varying separation.}
\label{tab:exp4}
\begin{tabular}{rlrrrrrrc}
\toprule
$s$ & Regime & Pred $R_A$ & Sim $R_A$ & Pred $R_B$ & Sim $R_B$ & Pred $R_C$ & Sim $R_C$ & A$=$C \\
\midrule
10 & all\_conf & 2.4611 & 2.4611 & 2.4611 & 2.4611 & 2.4611 & 2.4611 & \checkmark \\
35 & all\_conf & 2.4611 & 2.4611 & 2.4611 & 2.4611 & 2.4611 & 2.4611 & \checkmark \\
40 & mixed    & 1.8606 & 1.8606 & 2.4611 & 2.4611 & 1.8606 & 1.8606 & \checkmark \\
50 & mixed    & 2.1741 & 2.1741 & 2.4611 & 2.4611 & 2.1741 & 2.1741 & \checkmark \\
70 & mixed    & 2.8403 & 2.8403 & 2.7203 & 2.7203 & 2.8403 & 2.8403 & \checkmark \\
75 & all\_hid & 3.2250 & 3.2250 & 2.8667 & 2.8667 & 3.2250 & 3.2250 & \checkmark \\
100 & all\_hid & 4.3000 & 4.3000 & 3.8222 & 3.8222 & 4.3000 & 4.3000 & \checkmark \\
150 & all\_hid & 6.4500 & 6.4500 & 5.1600 & 5.1600 & 6.4500 & 6.4500 & \checkmark \\
\bottomrule
\end{tabular}
\end{table}

All 11 test points match. The outer links A and C always have identical rates (geometric symmetry), confirming symmetric recalculation. In the mixed regime at $s=70\,$m, an interesting crossover occurs where A/C become faster than B, causing B to finish last and transition to a solo phase.

%% ────────────────────────────────────────────────────────────────
\section{Experiment~5: Staggered Transfer Start}
\label{sec:exp5}

\subsection{Setup}
Two parallel 30\,m links at 100\,m separation (hidden terminals). Link~B starts transferring immediately, while Link~A's transfer is delayed by $\delta$ seconds (its predecessor task has a large compute cost). This tests the simulator's ability to dynamically recalculate in-flight transfers when interference conditions change mid-transfer.

\subsection{Predictions}
Let $R_0 = 8.6\MBs$ (solo rate) and $R_1 = 4.3\MBs$ (both-active SINR rate). Three phases occur:
\begin{enumerate}
    \item Phase~1: B transmits solo at $R_0$ for $\delta$ seconds.
    \item Phase~2: Both active at $R_1$ until B finishes.
    \item Phase~3: A transmits solo at $R_0$ until A finishes.
\end{enumerate}
By symmetry, both transfers have the same total duration: $T = \delta + (D - \delta R_0)/R_1$.

\subsection{Results}

\begin{table}[h]
\centering
\caption{Experiment~5: Staggered start with hidden terminals.}
\label{tab:exp5}
\begin{tabular}{rrrrrr}
\toprule
Delay Cost & $\delta$ (s) & Pred.\ Duration & Dur.\ A & Dur.\ B & A$=$B \\
\midrule
10000 & 0.1000 & 2.225591 & 2.225591 & 2.225591 & \checkmark \\
30000 & 0.3000 & 2.025591 & 2.025591 & 2.025591 & \checkmark \\
50000 & 0.5000 & 1.825591 & 1.825591 & 1.825591 & \checkmark \\
80000 & 0.8000 & 1.525591 & 1.525591 & 1.525591 & \checkmark \\
\bottomrule
\end{tabular}
\end{table}

All 4 test points match exactly. The A$=$B symmetry confirms that when A starts and B is recalculated, B's rate degradation exactly mirrors A's, and when B finishes, A correctly recovers to the solo rate.

%% ────────────────────────────────────────────────────────────────
\section{Experiment~6: Per-Link Bandwidth Sharing}
\label{sec:exp6}

\subsection{Setup}
Multiple concurrent flows (3 or 4) share a single 30\,m link with no inter-link interference. Each flow carries a different data size. Source tasks are queued on the same node (FIFO), creating a negligible ${\sim}10\,\mu$s stagger between transfer starts.

This experiment isolates the execution engine's per-link fair-sharing mechanism: $N$ flows on one link each receive $R_{\text{base}}/N$ of the bandwidth. When a flow completes, remaining flows are recalculated with a larger share.

\subsection{Multi-Phase Prediction}
With data sizes $D_1 < D_2 < \ldots < D_N$ and link bandwidth $B$:
\begin{itemize}    \item Phase~1: All $N$ flows at $B/N$. The smallest ($D_1$) finishes at $t_1 = D_1 \cdot N / B$.
    \item Phase~2: $N{-}1$ flows at $B/(N{-}1)$. The next-smallest finishes after its remaining data is transferred.
    \item Continue until the last flow completes solo at rate~$B$.
\end{itemize}

\subsection{Results}

\begin{table}[h]
\centering
\caption{Experiment~6: Per-link bandwidth sharing with multiple flows.}
\label{tab:exp6}
\begin{tabular}{llrrr}
\toprule
Test Case & Flow & Data (MB) & Pred.\ Duration (s) & Sim.\ Duration (s) \\
\midrule
3-flow & Ts0 &  5.0 & 1.744186 & 1.744161 \\
3-flow & Ts1 & 10.0 & 2.906977 & 2.906962 \\
3-flow & Ts2 & 15.0 & 3.488372 & 3.488352 \\
\midrule
3-flow-eq & Ts0 &  5.0 & 1.744186 & 1.744161 \\
3-flow-eq & Ts1 &  5.0 & 1.744186 & 1.744171 \\
3-flow-eq & Ts2 & 15.0 & 2.906977 & 2.906957 \\
\midrule
4-flow & Ts0 &  3.0 & 1.395349 & 1.395306 \\
4-flow & Ts1 &  6.0 & 2.441860 & 2.441837 \\
4-flow & Ts2 &  9.0 & 3.139535 & 3.139512 \\
4-flow & Ts3 & 12.0 & 3.488372 & 3.488342 \\
\bottomrule
\end{tabular}
\end{table}

All 10 test points match within tolerance. The ``3-flow-eq'' case validates that two flows with equal data sizes complete at essentially the same time (within microseconds), and the remaining flow correctly transitions to a higher rate. The ``4-flow'' case exercises four sequential phase transitions as flows complete one by one.

%% ────────────────────────────────────────────────────────────────
\section{Experiment~7: $N$-Way Bianchi Contention Scaling}
\label{sec:exp7}

\subsection{Setup}
$N$ parallel 30\,m links are placed at 5\,m vertical separation, all within carrier-sensing range of each other. This creates a complete conflict graph of size~$N$. Each link carries a 10\,MB transfer. Since all links have the same rate, they form a single-phase scenario.

\subsection{Predictions}
With all $N$ links in a complete conflict graph, the factor for each link is $\eta(N)/N$, giving $R = R_{\text{base}} \cdot \eta(N)/N = 8.6 \cdot \eta(N)/N$.

\subsection{Results}

\begin{table}[h]
\centering
\caption{Experiment~7: Bianchi contention scaling with $N$ links.}
\label{tab:exp7}
\begin{tabular}{rrrrrr}
\toprule
$N$ & $\eta(N)$ & $\eta(N)/N$ & Pred.\ Rate (MB/s) & Sim.\ Rate (MB/s) & Match \\
\midrule
2 & 0.8807 & 0.4404 & 3.7871 & 3.7871 & \checkmark \\
3 & 0.8585 & 0.2862 & 2.4611 & 2.4611 & \checkmark \\
4 & 0.8367 & 0.2092 & 1.7989 & 1.7989 & \checkmark \\
5 & 0.8179 & 0.1636 & 1.4067 & 1.4067 & \checkmark \\
6 & 0.8019 & 0.1336 & 1.1494 & 1.1494 & \checkmark \\
7 & 0.7882 & 0.1126 & 0.9683 & 0.9683 & \checkmark \\
8 & 0.7257 & 0.0907 & 0.7802 & 0.7802 & \checkmark \\
\bottomrule
\end{tabular}
\end{table}

All 7 test points match. The MAC efficiency $\eta(N)$ decreases gracefully from 0.88 at $N\!=\!2$ to 0.73 at $N\!=\!8$, while the per-link share $\eta(N)/N$ decreases much faster due to the $1/N$ division. All $N$ links within each test case receive identical rates, confirming symmetric contention behavior.

%% ────────────────────────────────────────────────────────────────
\section{Experiment~8: Five-Link Hidden-Terminal Cascade}
\label{sec:exp8}

\subsection{Setup}
Five parallel 30\,m links at 100\,m vertical separation (all pairs beyond $d_{\text{CS}}$, hence all hidden terminals). Data sizes are $\{2, 4, 6, 8, 10\}\,$MB to force links to complete at different times, creating up to four cascading phase transitions. The five links have asymmetric interference geometry: outer links (at $y=0$ and $y=400$) have one near and three increasingly distant interferers, while the middle link (at $y=200$) has two near and two far interferers.

\subsection{Multi-Phase Prediction}
A general cascade prediction is used:
\begin{enumerate}
    \item Compute SINR-based rates for all active links given current interferer set.
    \item Advance time until the next link completes.
    \item Remove completed link; recompute rates with fewer interferers.
    \item Repeat until all links finish.
\end{enumerate}
This produces up to 5 phases with successively fewer active links.

\subsection{Results}

\begin{table}[h]
\centering
\caption{Experiment~8: Five-link hidden-terminal cascade.}
\label{tab:exp8}
\begin{tabular}{rrrrrr}
\toprule
Link & Position $y$ & Data (MB) & Pred.\ Duration (s) & Sim.\ Duration (s) & Match \\
\midrule
$l_0$ &   0 &  2.0 & 0.465116 & 0.465116 & \checkmark \\
$l_1$ & 100 &  4.0 & 1.046512 & 1.046512 & \checkmark \\
$l_2$ & 200 &  6.0 & 1.656977 & 1.656977 & \checkmark \\
$l_3$ & 300 &  8.0 & 2.274709 & 2.274709 & \checkmark \\
$l_4$ & 400 & 10.0 & 2.300145 & 2.300145 & \checkmark \\
\bottomrule
\end{tabular}
\end{table}

All 5 test points match. This experiment is the most demanding stress test: five links undergo four phase transitions, each requiring correct \texttt{data\_remaining} tracking and SINR recomputation with a changing active-link set. Notably, links $l_3$ and $l_4$ finish at nearly the same time despite having different data sizes (8 vs.\ 10\,MB), because $l_4$ (an outer link) has a higher rate than $l_3$ (an inner link) in all phases.

%% ────────────────────────────────────────────────────────────────
\section{Experiment~9: Combined Bandwidth Sharing and Interference}
\label{sec:exp9}

\subsection{Setup}
This experiment combines both mechanisms: two flows share link~$l_{01}$ (n0$\to$n1, 30\,m) while a third flow on link~$l_{23}$ (n2$\to$n3, 30\,m at 100\,m separation) acts as a hidden terminal. The effective per-flow rate on $l_{01}$ is therefore:
\begin{equation}
    R_{\text{flow}} = \frac{R_{\text{base}} \cdot f_{\text{SINR}}}{N_{\text{flows}}} = \frac{8.6 \times 0.5}{2} = 2.15 \MBs
\end{equation}
where $f_{\text{SINR}} = R_{\text{SINR}}(30, 104.4) / R_{\text{base}} = 4.3/8.6 = 0.5$ accounts for the hidden terminal, and $N_{\text{flows}} = 2$ accounts for per-link fair sharing. Three test cases exercise different phase orderings by varying data sizes.

\subsection{Phase Analysis (Case~1: $l_{01}$: 5/10\,MB, $l_{23}$: 8\,MB)}
\begin{itemize}    \item \textbf{Phase~1} (all active): $l_{01}$ flows at 2.15\,MB/s each; $l_{23}$ at 4.3\,MB/s. The $l_{23}$ flow (8\,MB) finishes first at $8/4.3 = 1.86\,$s.
    \item \textbf{Phase~2} ($l_{01}$ only, 2 flows, no interference): Each flow at $8.6/2 = 4.3\MBs$. The 5\,MB flow finishes.
    \item \textbf{Phase~3} ($l_{01}$ only, 1 flow): Solo at $8.6\MBs$. The 10\,MB flow finishes.
\end{itemize}

\subsection{Results}

\begin{table}[h]
\centering
\caption{Experiment~9: Combined bandwidth sharing and hidden-terminal interference.}
\label{tab:exp9}
\begin{tabular}{llrrrr}
\toprule
Case & Flow & Data (MB) & Link & Pred.\ Duration (s) & Sim.\ Duration (s) \\
\midrule
1 & Ta0 &  5.0 & $l_{01}$ & 2.093023 & 2.093018 \\
1 & Tb0 & 10.0 & $l_{01}$ & 2.674419 & 2.674408 \\
1 & Tc0 &  8.0 & $l_{23}$ & 1.860465 & 1.860465 \\
\midrule
2 & Ta0 &  5.0 & $l_{01}$ & 1.511628 & 1.511623 \\
2 & Tb0 & 10.0 & $l_{01}$ & 2.093023 & 2.093013 \\
2 & Tc0 &  3.0 & $l_{23}$ & 0.697674 & 0.697674 \\
\midrule
3 & Ta0 &  3.0 & $l_{01}$ & 1.395349 & 1.395339 \\
3 & Tb0 & 15.0 & $l_{01}$ & 4.186047 & 4.186037 \\
3 & Tc0 & 20.0 & $l_{23}$ & 4.418605 & 4.418605 \\
\bottomrule
\end{tabular}
\end{table}

All 9 test points match. Case~1 has the $l_{23}$ flow finishing first, removing interference before the $l_{01}$ sharing phases complete. Case~2 has $l_{23}$ finishing very early, so most of the $l_{01}$ transfer occurs without interference. Case~3 has the first $l_{01}$ flow finishing before $l_{23}$, testing the case where per-link sharing changes while interference is still active.

%% ────────────────────────────────────────────────────────────────
\section{Summary and Conclusions}
\label{sec:conclusions}

\subsection{Results Summary}

\begin{table}[h]
\centering
\caption{Summary of all verification experiments.}
\label{tab:summary}
\begin{tabular}{rlrc}
\toprule
\# & Experiment & Test Points & Result \\
\midrule
1 & Link Length vs.\ Data Rate & 20 & PASS \\
2 & Parallel Link Separation & 15 & PASS \\
3 & Two Transmitters to Same Receiver & 18 & PASS \\
4 & Three Parallel Links (Multi-Phase) & 11 & PASS \\
5 & Staggered Transfer Start & 4 & PASS \\
6 & Per-Link Bandwidth Sharing & 10 & PASS \\
7 & $N$-Way Bianchi Scaling ($N=2\ldots 8$) & 7 & PASS \\
8 & Five-Link Hidden-Terminal Cascade & 5 & PASS \\
9 & Combined Sharing + Interference & 9 & PASS \\
\midrule
& \textbf{Total} & \textbf{99} & \textbf{ALL PASS} \\
\bottomrule
\end{tabular}
\end{table}

\subsection{Conclusions}

All 99 individual test points across 9 experiments match analytical predictions within the 1\% tolerance. The experiments verify the following aspects of the \ncsim{} interference model and execution engine:

\begin{enumerate}    \item \textbf{PHY-layer baseline} (Exp.~1): The SNR-to-MCS-to-rate pipeline produces correct rates across the full distance range, including the out-of-range fallback.

    \item \textbf{Conflict-graph contention} (Exps.~2, 3, 4, 7): The Bianchi MAC efficiency formula is correctly applied for $N = 2$ through $N = 8$ contending stations, with symmetric rate assignment.

    \item \textbf{Hidden-terminal SINR degradation} (Exps.~2, 3, 4, 8): SINR computation correctly aggregates interference from multiple hidden terminals, with proper MCS re-selection and minimum-factor clamping.

    \item \textbf{Mixed contention + hidden-terminal regimes} (Exp.~4): The model correctly composes Bianchi contention sharing with SINR degradation when some links are in the conflict graph and others are not.

    \item \textbf{Dynamic recalculation} (Exps.~4, 5, 8): When transfers start or complete, all affected links are recalculated with correct partial-progress tracking (\texttt{data\_remaining}), even through 4+ cascading phase transitions.

    \item \textbf{Per-link fair sharing} (Exps.~6, 9): Multiple flows on the same physical link correctly receive equal bandwidth shares, with proper recalculation when flows complete.

    \item \textbf{Mechanism composition} (Exp.~9): Per-link bandwidth sharing and inter-link hidden-terminal interference compose correctly, with phase transitions triggered by either mechanism.
\end{enumerate}

These results provide confidence that the \texttt{csma\_bianchi} interference model in \ncsim{} faithfully implements the intended analytical model and that the execution engine's dynamic recalculation machinery is correct.

%% ────────────────────────────────────────────────────────────────
\begin{thebibliography}{9}
\bibitem{bianchi2000}
G.~Bianchi, ``Performance Analysis of the IEEE 802.11 Distributed Coordination Function,'' \emph{IEEE Journal on Selected Areas in Communications}, vol.~18, no.~3, pp.~535--547, Mar.~2000.
\end{thebibliography}

\end{document}
